% ============================================================
% Textbook Authoring System — Color Style Template
% ============================================================

\usepackage{xcolor}
\usepackage{titlesec}
\usepackage{mdframed}
\usepackage{booktabs}
\usepackage{colortbl}
\usepackage{fancyhdr}
\usepackage{tcolorbox}
\usepackage{geometry}
\usepackage{hyperref}
\usepackage{enumitem}

\tcbuselibrary{skins, breakable}

% ── 색상 팔레트 ──────────────────────────────────────────────
\definecolor{primary}{HTML}{1A3A5C}      % 진한 네이비 — H1
\definecolor{secondary}{HTML}{2E7BAE}    % 미드 블루  — H2
\definecolor{accent}{HTML}{3AAFA9}       % 틸       — H3
\definecolor{highlight}{HTML}{F0F7FF}    % 연한 하늘 — 배경
\definecolor{tableheader}{HTML}{1A3A5C}  % 테이블 헤더
\definecolor{tablerowodd}{HTML}{F7FBFF}  % 테이블 홀수행
\definecolor{tableroweven}{HTML}{FFFFFF} % 테이블 짝수행
\definecolor{rulecolor}{HTML}{2E7BAE}    % 구분선
\definecolor{boxbg}{HTML}{EFF9F8}        % 박스 배경 (틸 연하게)
\definecolor{boxborder}{HTML}{3AAFA9}    % 박스 테두리
\definecolor{warningbg}{HTML}{FFF8E7}    % 경고 박스 배경
\definecolor{warningborder}{HTML}{F0A500}% 경고 박스 테두리
\definecolor{codebg}{HTML}{F4F6F8}       % 코드 블록 배경
\definecolor{linkcolor}{HTML}{2E7BAE}    % 링크 색상

% ── 페이지 여백 ──────────────────────────────────────────────
\geometry{
  top=2.5cm, bottom=2.5cm,
  left=2.8cm, right=2.8cm,
  headheight=14pt
}

% ── 헤딩 스타일 (계층별 색상) ────────────────────────────────
% H1: 진한 네이비, 굵게, 하단 구분선
\titleformat{\section}
  {\color{primary}\Large\bfseries}
  {\color{primary}\thesection}{1em}{}
  [\color{rulecolor}\titlerule]

% H2: 미드 블루, 굵게
\titleformat{\subsection}
  {\color{secondary}\large\bfseries}
  {\color{secondary}\thesubsection}{1em}{}

% H3: 틸, 보통 굵기
\titleformat{\subsubsection}
  {\color{accent}\normalsize\bfseries}
  {\color{accent}\thesubsubsection}{1em}{}

% ── 헤더/푸터 ────────────────────────────────────────────────
\pagestyle{fancy}
\fancyhf{}
\fancyhead[L]{\small\color{secondary}\leftmark}
\fancyhead[R]{\small\color{secondary}Textbook Authoring System}
\fancyfoot[C]{\small\color{secondary}\thepage}
\renewcommand{\headrulewidth}{0.4pt}
\renewcommand{\headrule}{\color{rulecolor}\hrule width\headwidth height\headrulewidth}

% ── 테이블 스타일 ────────────────────────────────────────────
\usepackage{array}
\renewcommand{\arraystretch}{1.4}

% 테이블 헤더 행 자동 색상 (longtable/tabular 공통)
\usepackage{longtable}
\AtBeginEnvironment{longtable}{%
  \rowcolors{2}{tablerowodd}{tableroweven}
}
\AtBeginEnvironment{tabular}{%
  \rowcolors{2}{tablerowodd}{tableroweven}
}

% ── 코드 블록 ────────────────────────────────────────────────
\usepackage{fancyvrb}
\usepackage{framed}

% Shaded는 pandoc이 코드블록 있을 때만 정의하므로 조건부 처리
\makeatletter
\AtBeginDocument{%
  \@ifundefined{Shaded}{%
    \newenvironment{Shaded}{%
      \begin{tcolorbox}[enhanced,colback=codebg,colframe=secondary!30,
        boxrule=0.5pt,arc=4pt,left=6pt,right=6pt,top=4pt,bottom=4pt,breakable]
    }{%
      \end{tcolorbox}
    }%
  }{%
    \renewenvironment{Shaded}{%
      \begin{tcolorbox}[enhanced,colback=codebg,colframe=secondary!30,
        boxrule=0.5pt,arc=4pt,left=6pt,right=6pt,top=4pt,bottom=4pt,breakable]
    }{%
      \end{tcolorbox}
    }%
  }%
}
\makeatother

% ── 강조 박스 (blockquote → 틸 박스) ─────────────────────────
\renewenvironment{quote}{%
  \begin{tcolorbox}[
    enhanced,
    colback=boxbg,
    colframe=boxborder,
    boxrule=1.5pt,
    leftrule=4pt,
    arc=3pt,
    left=8pt, right=6pt, top=4pt, bottom=4pt,
    breakable
  ]
  \itshape
}{%
  \end{tcolorbox}
}

% ── 학습목표 박스 (문서 첫 번째 구분용) ─────────────────────
\newenvironment{objectivebox}{%
  \begin{tcolorbox}[
    enhanced,
    title={\bfseries 학습 목표},
    colbacktitle=primary,
    colframe=primary,
    colback=highlight,
    fonttitle=\color{white}\bfseries,
    boxrule=1pt,
    arc=4pt,
    breakable
  ]
}{%
  \end{tcolorbox}
}

% ── 하이퍼링크 색상 ──────────────────────────────────────────
\hypersetup{
  colorlinks=true,
  linkcolor=secondary,
  urlcolor=accent,
  citecolor=secondary,
  pdfborder={0 0 0}
}

% ── 수식 번호 색상 ───────────────────────────────────────────
\makeatletter
\renewcommand{\theequation}{\color{secondary}\arabic{equation}}
\makeatother

% ── 목차 스타일 ──────────────────────────────────────────────
\usepackage{tocloft}
\renewcommand{\cftsecfont}{\color{primary}\bfseries}
\renewcommand{\cftsubsecfont}{\color{secondary}}
\renewcommand{\cftsubsubsecfont}{\color{accent}}
\renewcommand{\cftsecpagefont}{\color{primary}\bfseries}
